\begin{textsample}{POS Dim 3 – pi – Score 7.00 – pi/pi\_inf\_3.txt}  \label{ex:f3_pos_266}
2.
MARCOS LEGAIS A política da educação \textbf{infantil} ganhou força que sustenta a sua implementação a partir dos anos 1990.
Nesse período foi promulgada a Lei de Diretrizes e Base da Educação – LDB ( Lei nº 9.394/1996 ).
Essa Lei apresenta a necessidade de um currículo que traz a sua composição por área, faixa etária e eixo.
Art. 29 – A educação \textbf{infantil}, primeira etapa da educação básica, tem como finalidade o desenvolvimento integral da \textbf{criança} até os seis anos de idade, em seus aspectos físico, psicológico, intelectual e social, complementando a ação da família e da comunidade.
Art. 30 – A educação \textbf{infantil} será oferecida em : I – \textbf{creches} ou entidades equivalentes, para \textbf{crianças} de até três anos de idade ; II – \textbf{pré-escolas} para \textbf{crianças} de quatro a seis anos de idade.
Art. 31 – Na educação \textbf{infantil} a avaliação far -se-á mediante \textbf{acompanhamento} e \textbf{registro} de seu desenvolvimento, sem o objetivo de promoção, mesmo para o acesso ao ensino fundamental ( Lei nº 9.394/1996 ).
A educação integral já é prerrogativa na LDB, pois reforça nossa intencionalidade de implementar um currículo com foco em uma educação plena para as escolas piauienses.
Dessa forma, fica evidente que a educação \textbf{infantil} é uma oportunidade para o pleno desenvolvimento do educando.
A legislação supracitada define também a educação \textbf{infantil} como a primeira etapa da educação básica.
Assim fica assegurado que a \textbf{criança} de 0 a 5 anos e 11 \textbf{meses} de idade tem direito a ser matriculada em \textbf{creches} e \textbf{pré-escolas}, ficando claro o compromisso da educação \textbf{infantil} em desenvolver uma educação de qualidade de forma plena.
Na década de 2000 foi fortalecida essa legislação educacional para a educação \textbf{infantil} com a implementação das Diretrizes Curriculares Nacionais para a Educação \textbf{Infantil}, DCNEI, que reforça a articulação da EI como etapa integrante da educação básica e define um currículo para essa etapa educacional.
O currículo da Educação \textbf{Infantil} é concebido como um conjunto de práticas que buscam articular as experiências e os saberes das \textbf{crianças} com os conhecimentos que fazem parte do patrimônio cultural, artístico, ambiental, científico e tecnológico, de modo a promover o desenvolvimento integral de \textbf{crianças} de 0 a 5 anos de idade ( BRASIL, 2009, Art. 3º ).
A contextualização do processo educacional é um aspecto que fica evidente nessa concepção de currículo proposta pelas DCNEI que é reforçada pelo Referencial Curricular Nacional para Educação \textbf{Infantil} – RCNEI, ao defender a implementação de propostas pedagógicas compatíveis com as especificidades de cada região do país.
Essa legislação é compatível com o que determina a Constituição Federal de 1988 e o Estatuto da Criança e do Adolescente – ECA ( 1990 ), sobre os direitos educacionais da \textbf{criança} e do adolescente.

% matched lemmas: acompanhamento, creche, criança, infantil, mês, pré-escola, registro
\end{textsample}
